\documentclass[11pt]{article}

% ---------------------------------------------------------------------------
% Packages
% ---------------------------------------------------------------------------
\usepackage[utf8]{inputenc}
\usepackage[T1]{fontenc}
\usepackage{lmodern}
\usepackage[margin=1in]{geometry}
\usepackage{graphicx}
\usepackage{amsmath}
\usepackage{booktabs}
\usepackage{xcolor}
\usepackage{listings}
\usepackage{authblk}
\usepackage[numbers,sort&compress]{natbib}
\usepackage{hyperref}
\usepackage{cleveref}
\usepackage{lineno}
\linenumbers

\hypersetup{colorlinks=true, linkcolor=blue!60!black, citecolor=blue!60!black, urlcolor=blue!60!black}

% Inline notes for drafting
\newcommand{\todo}[1]{\textcolor{red}{\textbf{[TODO: #1]}}}

% Placeholder figure box until real figures exist
\newcommand{\placeholderfig}[2]{%
  \fbox{\parbox[c][#2][c]{0.95\linewidth}{\centering\textit{#1}}}}

\lstset{
  language=Python,
  basicstyle=\ttfamily\small,
  keywordstyle=\color{blue!60!black}\bfseries,
  commentstyle=\color{gray}\itshape,
  stringstyle=\color{green!40!black},
  frame=single,
  breaklines=true,
  columns=fullflexible,
  showstringspaces=false,
}

\newcommand{\pynaviz}{\texttt{pynaviz}}
\newcommand{\pynapple}{\texttt{pynapple}}

% ---------------------------------------------------------------------------
% Title / authors
% ---------------------------------------------------------------------------
\title{Pynaviz: GPU-accelerated, synchronized visualization of neural and behavioral time series in Python}

\author[1]{Edoardo Balzani}
\author[1]{Guillaume Viejo}
\author{\todo{other authors and order}}
\affil[1]{\todo{Affiliation}}

\date{\today}

\begin{document}

\maketitle

% ---------------------------------------------------------------------------
\begin{abstract}
Modern systems neuroscience experiments routinely combine high-channel-count
electrophysiology or imaging with high-frame-rate behavioral video and
pose-estimation outputs. Exploring these heterogeneous streams together, at
their native resolution and on a common time axis, remains a practical
bottleneck: general-purpose plotting libraries are not designed for
interactive navigation of hours-long recordings, while dedicated viewers are
often tied to a single data modality or file format. Here we present
\pynaviz{}, an open-source Python package for interactive, GPU-accelerated
visualization of neurophysiological and behavioral data. \pynaviz{} is built
on the \pynapple{} data model \citep{viejo2023pynapple} and on the
WebGPU-based rendering engine \texttt{pygfx}, and provides one visual per
\pynapple{} object (\texttt{Tsd}, \texttt{TsdFrame}, \texttt{TsdTensor},
\texttt{TsGroup}, \texttt{Ts}, \texttt{IntervalSet}) as well as
frame-accurate video playback. Any number of views can be synchronized in
time, so that panning in a raw LFP trace, a spike raster or a pose trajectory
moves the corresponding video frame and vice versa. Large datasets are
streamed from disk on demand, metadata can be used to sort, group and color
signals interactively, and complete workspaces can be saved and restored. \pynaviz{}
can be used from a script or notebook, as a standalone desktop application
launched from the command line on NWB, video or raw electrophysiology files,
or embedded in custom Qt applications.
\todo{one sentence on validation / benchmark results.}
\end{abstract}

% ---------------------------------------------------------------------------
\section{Introduction}
\label{sec:intro}

\todo{Motivation paragraph.} Visual inspection of raw data is a necessary
step in every neurophysiology pipeline: checking spike sorting against the
raw signal, verifying the alignment of behavioral events, identifying
artifacts, or simply building intuition before quantitative analysis.
Recordings now routinely span hundreds to thousands of channels
\citep{jun2017neuropixels,steinmetz2021neuropixels}, are paired with one or
several high-speed cameras, and are augmented by markerless pose-estimation
\citep{mathis2018deeplabcut,pereira2022sleap}. The resulting datasets are
heterogeneous in sampling rate, dimensionality and storage format.

\todo{Gap paragraph.} Static plotting libraries such as Matplotlib
\citep{hunter2007matplotlib} are ill-suited for interactive browsing of long
recordings. Several interactive tools exist (e.g.\ Neuroscope
\citep{hazan2006klusters}, ephyviewer, napari \citep{napari}, fastplotlib,
SpikeInterface GUIs \citep{buccino2020spikeinterface}) but are generally
tied to a specific data type (continuous signals, images, or spikes) or a
specific file format, and do not offer native, frame-accurate synchronization
between neural signals and behavioral video. \todo{Refine comparison; add
table?}

\todo{Contribution paragraph.} We introduce \pynaviz{}, which closes this gap
by (i) mapping each core \pynapple{} object to a dedicated GPU-rendered
visual, (ii) synchronizing an arbitrary number of visuals on a shared time
axis, including video, and (iii) exposing the same functionality through a
Python API, a Qt desktop application and a command-line entry point.

% ---------------------------------------------------------------------------
\section{Design and implementation}
\label{sec:design}

\subsection{Data model}
\pynaviz{} adopts the data model of \pynapple{} \citep{viejo2023pynapple}:
timestamps (\texttt{Ts}), univariate (\texttt{Tsd}), multivariate
(\texttt{TsdFrame}) and tensor-valued (\texttt{TsdTensor}) time series,
collections of point processes (\texttt{TsGroup}) and sets of epochs
(\texttt{IntervalSet}). Metadata attached to these objects (e.g.\ brain
region, cell type, channel depth) is directly accessible to the visuals.
Because \pynapple{} reads NWB files \citep{rubel2022nwb} and, through Neo
\citep{garcia2014neo}, most raw electrophysiology formats, \pynaviz{}
inherits this I/O layer without additional code.

\Cref{tab:visuals} lists the visual associated with each object.

\begin{table}[h]
\centering
\caption{Mapping between \pynapple{} objects and \pynaviz{} visuals.}
\label{tab:visuals}
\begin{tabular}{@{}lll@{}}
\toprule
Object & Visual & Display modes \\
\midrule
\texttt{Tsd}         & \texttt{PlotTsd}         & line \\
\texttt{TsdFrame}    & \texttt{PlotTsdFrame}    & stacked lines, heat map, $x$ vs.\ $y$ \\
\texttt{TsdTensor}   & \texttt{PlotTsdTensor}   & image sequence \\
\texttt{TsGroup}     & \texttt{PlotTsGroup}     & raster \\
\texttt{Ts}          & \texttt{PlotTs}          & event markers \\
\texttt{IntervalSet} & \texttt{PlotIntervalSet} & shaded epochs / overlay \\
Video file           & \texttt{PlotVideo}       & frames (+ skeleton overlay) \\
\bottomrule
\end{tabular}
\end{table}

\subsection{GPU rendering}
All visuals are rendered with \texttt{pygfx}, a Python rendering engine built
on WebGPU (\texttt{wgpu-py}). Data are uploaded once to GPU buffers and
navigation (pan, zoom) is implemented as camera transformations, so that
interaction cost is largely independent of the number of samples on screen.
\todo{Describe buffer layout for TsdFrame / TsGroup, line vs.\ image modes.}

\subsection{Time synchronization}
Each visual owns a \emph{controller} that translates user input into camera
updates. Controllers are registered to a \texttt{ControllerGroup}, which
broadcasts synchronization events (pan, zoom, go-to-time) to all other
members. Time-series visuals use span controllers that preserve the visible
time window, whereas video and tensor visuals use a frame controller that
maps the current time to the nearest frame. Synchronization rules are
modular, allowing views with different units or aspect ratios to follow the
same time axis. \todo{Figure: architecture diagram.}

\subsection{Streaming large datasets}
For memory-mapped or lazily loaded continuous data (e.g.\ raw wide-band
recordings), \pynaviz{} streams a fixed-size window around the current view
position rather than loading the full array. The window is refreshed as the
user pans or zooms, keeping GPU memory bounded. \todo{Numbers.}

\subsection{Video decoding}
Video is decoded with PyAV (FFmpeg bindings). A background thread builds an
index of keyframe presentation timestamps at load time, enabling
frame-accurate random seeks, while a frame buffer amortizes sequential
playback. Video timestamps can be provided explicitly, so that
camera frames align with the neural clock. Pose-estimation outputs (e.g.\
DeepLabCut, SLEAP) can be overlaid as skeletons on top of the video.
\todo{Seek latency benchmark.}

\subsection{Metadata-driven interaction}
Visuals expose \texttt{sort\_by}, \texttt{group\_by} and \texttt{color\_by}
operations driven by object metadata, available both programmatically and
from the GUI. \texttt{IntervalSet} objects can be superimposed on any
time-series visual and used to jump between epochs (e.g.\ trials, sleep
states).

\subsection{User interfaces}
\pynaviz{} can be used at three levels:
\begin{enumerate}
  \item \textbf{Python API.} Individual visuals (\texttt{PlotTsd}, \ldots)
        rendered in a standalone window or a Jupyter notebook.
  \item \textbf{Qt application.} The \texttt{scope} function opens a
        dockable multi-panel window in which variables from a dictionary, a
        namespace, an NWB file or a raw recording can be added, arranged and
        synchronized. Layouts are serialized to JSON and can be restored.
  \item \textbf{Command line.} \texttt{pynaviz [files] [-l layout.json]}
        opens the same application directly on NWB, NPZ, video or raw
        electrophysiology files.
\end{enumerate}

\begin{lstlisting}[caption={Minimal example.},label={lst:example}]
import pynapple as nap
from pynaviz import scope

data = nap.load_file("session.nwb")
scope({"lfp": data["lfp"], "units": data["units"],
       "video": "camera.mp4"})
\end{lstlisting}

\begin{figure}[t]
\centering
\placeholderfig{Figure 1: overview of the \pynaviz{} Qt application with
synchronized LFP, spike raster, epochs and video.}{6cm}
\caption{\todo{Overview figure.}}
\label{fig:overview}
\end{figure}

% ---------------------------------------------------------------------------
\section{Use cases}
\label{sec:usecases}

\subsection{Head-direction cells and behavior}
\todo{Spike raster sorted by preferred direction, synchronized with head
angle and tuning curves. Based on \texttt{docs/examples/example\_head\_direction.py}.}

\subsection{Raw LFP and sleep states}
\todo{Wide-band streaming, IntervalSet overlay of sleep/wake epochs. Based on
\texttt{example\_lfp.py}.}

\subsection{Pose estimation and video}
\todo{DeepLabCut output overlaid on video, $x$ vs.\ $y$ trajectory view.
Based on \texttt{example\_dlc\_pose.py}.}

\subsection{Trial-based experiments}
\todo{Jumping between trials, multi-camera playback. Based on
\texttt{example\_trials.py} and \texttt{example\_videos.py}.}

\begin{figure}[t]
\centering
\placeholderfig{Figure 2: use-case panels A--D.}{8cm}
\caption{\todo{Use cases.}}
\label{fig:usecases}
\end{figure}

% ---------------------------------------------------------------------------
\section{Performance}
\label{sec:performance}
\todo{Benchmarks: frame rate vs.\ number of channels / samples on screen;
video seek latency; memory footprint with streaming; comparison with
Matplotlib / other viewers. Hardware description.}

\begin{figure}[t]
\centering
\placeholderfig{Figure 3: benchmarks.}{5cm}
\caption{\todo{Performance.}}
\label{fig:benchmarks}
\end{figure}

% ---------------------------------------------------------------------------
\section{Discussion}
\label{sec:discussion}
\todo{Summary; limitations (e.g.\ GPU requirement, Qt dependency for full
GUI); relation to the \pynapple{} ecosystem; future work (e.g.\ web/remote
rendering, 3D visuals, plugin system, tighter integration with analysis
outputs such as decoding).}

% ---------------------------------------------------------------------------
\section*{Code availability}
\pynaviz{} is open source and available at
\url{https://github.com/pynapple-org/pynaviz}, installable via
\texttt{pip install pynaviz[qt]}. Documentation is hosted at
\url{https://pynaviz.readthedocs.io}. \todo{License, version used, Zenodo DOI.}

\section*{Acknowledgments}
\todo{Funding, pygfx developers.}

\section*{Author contributions}
\todo{CRediT.}

\bibliographystyle{unsrtnat}
\bibliography{references}

\end{document}
